\begin{bankssolutions}{Solutions to Implicit Exercises for Chapter 5}

\BanksImplicitSolution{I-CH05-001}
The dual definition of $R_i$ is inverted by
\begin{equation*}
  J_{ij}=\epsilon_{ijk}R_k.
\end{equation*}
For example, the Lorentz algebra gives
\begin{equation*}
  [K_i,K_j]=[J_{0i},J_{0j}]
  =-\ii J_{ij}=-\ii\epsilon_{ijk}R_k.
\end{equation*}
The boost-rotation bracket follows first in the spatial-plane basis:
\begin{equation*}
  [K_i,J_{mn}]
  =-\ii\left(\delta_{im}K_n-\delta_{in}K_m\right).
\end{equation*}
Antisymmetry of the commutator and contraction with
$R_i=\epsilon_{imn}J_{mn}/2$ then give
\begin{equation*}
  [R_i,K_j]
  =\frac12\epsilon_{imn}[J_{mn},K_j]
  =+\ii\epsilon_{ijk}K_k.
\end{equation*}
For the rotation bracket, the definitions say
$R_1=J_{23}$, $R_2=J_{31}$, and $R_3=J_{12}$.  One representative
calculation is
\begin{equation*}
  [R_1,R_2]=[J_{23},J_{31}]
  =-\ii J_{21}=+\ii R_3.
\end{equation*}
Cyclic permutations supply the remaining components.  The resulting algebra
is
\begin{align*}
  [R_i,R_j]&=+\ii\epsilon_{ijk}R_k,\\
  [R_i,K_j]&=+\ii\epsilon_{ijk}K_k,\\
  [K_i,K_j]&=-\ii\epsilon_{ijk}R_k.
\end{align*}
The sign in the first line follows from Banks's choice
$J_{ij}=+\epsilon_{ijk}R_k$.  Reversing the orientation of all $R_i$ gives
the other common convention.

The independent-pair convention stated in the exercise is essential here.
If $\epsilon_{ijk}J_{ij}$ is evaluated as a sum over all ordered values of
$i,j$, it equals $2R_k$, so its relative normalization with $J_{0k}$ differs
from the displayed source combination.  In the calculation below,
$C_k^\pm$ denotes Banks's combination with each antisymmetric spatial pair
counted once.  The final rescaling to $T_k^\pm$ fixes the standard
$\mathfrak{su}(2)$ normalization.

Now form $C_i^\pm=R_i\pm\ii K_i$.  Substitution of the component relations
gives
\begin{align*}
  [C_i^\pm,C_j^\pm]
  &=[R_i,R_j]
    \mathbin{\pm}\ii[R_i,K_j]
    \mathbin{\pm}\ii[K_i,R_j]
    -[K_i,K_j]\\
  &=+2\ii\epsilon_{ijk}C_k^\pm,\\
  [C_i^+,C_j^-]&=0.
\end{align*}
The cancellation in the mixed bracket can be seen term by term:
\begin{equation*}
  [R_i,R_j]-\ii[R_i,K_j]
  +\ii[K_i,R_j]+[K_i,K_j]=0.
\end{equation*}
Hence
\begin{equation*}
  T_i^\pm=+\frac12C_i^\pm
\end{equation*}
obey
$[T_i^\pm,T_j^\pm]=\ii\epsilon_{ijk}T_k^\pm$ and
$[T_i^+,T_j^-]=0$.  These are the two copies of
$\mathfrak{su}(2)$ in the complexified Lorentz algebra.  Their sum recovers
spatial rotations:
\begin{equation*}
  R_i=T_i^++T_i^-,
  \qquad
  K_i=-\ii\left(T_i^+-T_i^-\right).
\end{equation*}
As a direct representation check, choose
$T_i^+=0$ and $T_i^-=\sigma_i/2$.  The reconstruction formula gives
$R_i=\sigma_i/2$ and $K_i=+\ii\sigma_i/2$, which satisfy all three Lorentz
brackets above.  Only one chiral algebra acts.  Exchanging the two $T$ sets
reverses the sign of the boost action and gives the conjugate chirality.

An irreducible finite-dimensional Lorentz representation may therefore be
labeled by one spin for each copy.  This is the algebraic origin of Banks's
$[n_L,n_R]$ labels and of the exchange of the two chiralities under spatial
reflection.

\BanksImplicitSolution{I-CH05-002}
Define
\begin{equation*}
  \sigma^{\mu\nu}=\frac\ii2[\gamma^\mu,\gamma^\nu].
\end{equation*}
The stated Hermiticity relation for the Dirac matrices implies
\begin{equation*}
  (\sigma^{\mu\nu})^\dagger
  =\gamma^0\sigma^{\mu\nu}\gamma^0.
\end{equation*}
For real Lorentz parameters one may write the spin representation as
\begin{equation*}
  S(\Lambda)=
  \exp\left(-\frac\ii4\omega_{\mu\nu}\sigma^{\mu\nu}\right).
\end{equation*}
Taking its Hermitian adjoint then gives
\begin{align*}
  S(\Lambda)^\dagger
  &=\gamma^0
    \exp\left(+\frac\ii4\omega_{\mu\nu}\sigma^{\mu\nu}\right)
    \gamma^0\\
  &=\gamma^0S(\Lambda)^{-1}\gamma^0.
\end{align*}
Multiplication on the right by $\gamma^0$ proves the required
pseudo-unitarity identity.

The transformed Dirac adjoint is consequently
\begin{equation*}
  \bar\psi'(x')
  =\psi(x)^\dagger S(\Lambda)^\dagger\gamma^0
  =\bar\psi(x)S(\Lambda)^{-1}.
\end{equation*}
The mass bilinear is a scalar:
\begin{equation*}
  \bar\psi'(x')\psi'(x')=\bar\psi(x)\psi(x).
\end{equation*}
For the derivative term, use
$\partial'_\mu=(\Lambda^{-1})^\nu{}_{\mu}\partial_\nu$ and the spinor
intertwining identity.  One obtains
\begin{align*}
  -\ii\bar\psi'\gamma^\mu\partial'_\mu\psi'
  &=\bar\psi S^{-1}(-\ii)\gamma^\mu S
    (\Lambda^{-1})^\nu{}_{\mu}\partial_\nu\psi\\
  &=-\ii\bar\psi\gamma^\nu\partial_\nu\psi.
\end{align*}
Both terms in the Dirac Lagrangian are Lorentz scalars.

The role of $\gamma^0$ is visible in a boost of rapidity $\chi$ along the
$x^1$ direction.  Up to the sign convention for $\chi$, its spin matrix is
\begin{equation*}
  S_B(\chi)=\exp\left(\frac\chi2\gamma^0\gamma^1\right).
\end{equation*}
The matrix $\gamma^0\gamma^1$ is Hermitian, so
$S_B^\dagger=S_B$ and $S_B^\dagger S_B=S_B^2$ depends on $\chi$.  Thus
$\psi^\dagger\psi$ changes under a boost.  Since
$\{\gamma^0,\gamma^0\gamma^1\}=0$, one instead has
\begin{equation*}
  S_B^\dagger\gamma^0S_B=\gamma^0.
\end{equation*}
The matrix $\gamma^0$ supplies the invariant intertwiner between the Dirac
representation and its Hermitian-conjugate representation, which is the
property used in Banks's definition of $\bar\psi$.

\BanksImplicitSolution{I-CH05-003}
Write
$B_{ij}=\epsilon_{ab}\psi_i^a\psi_j^b$.  Fermion anticommutation and the
antisymmetry of $\epsilon_{ab}$ make $B_{ij}=B_{ji}$, which explains why each
$Y^A$ is symmetric.  Under the stated CP action,
\begin{equation*}
  \phi_A B_{ij}
  \longmapsto
  Z_{AB}X_{ik}X_{jl}\phi_B^*B_{kl}^\dagger.
\end{equation*}
Comparison with the Hermitian-conjugate term in $\mathcal L_Y$ gives
\begin{equation*}
  (Y^B)^*
  =\sum_A Z_{AB}X^T Y^A X.
\end{equation*}
This equation, together with the corresponding conditions on the scalar
potential, is the generalized-CP invariance condition.

The involution condition converts the generalized CP action into an ordinary
real structure.  Indeed, unitarity and $XX^*=1$ imply
$X^*=X^{-1}=X^\dagger$, hence $X=X^T$.  The same argument gives $Z=Z^T$.
The Takagi factorization of these symmetric unitary matrices can be chosen as
\begin{equation*}
  X=VV^T,
  \qquad
  Z=WW^T,
\end{equation*}
with $V$ and $W$ unitary.  The compatible basis changes
\begin{equation*}
  \psi'=V^\dagger\psi,
  \qquad
  \phi'=W^\dagger\phi
\end{equation*}
turn the CP action into componentwise complex conjugation.  For instance,
\begin{equation*}
  \psi'\longmapsto
  V^\dagger X V^*\psi'^*=\psi'^*,
\end{equation*}
because $V^TV^*=1$.  The scalar transformation works in the same way.

In this CP basis the Yukawa matrices are
\begin{equation*}
  Y'^B=\sum_A W_{AB}V^T Y^A V.
\end{equation*}
CP sends each holomorphic Yukawa monomial to its Hermitian conjugate.  Equality
of their coefficients requires
\begin{equation*}
  Y'^B=(Y'^B)^*
\end{equation*}
for every $B$.  CP invariance therefore supplies a compatible basis in which
the full Yukawa collection is real.  Conversely, if such a basis exists, the
canonical CP transformation by complex conjugation preserves every Yukawa
term.  Undoing the basis change produces the matrices $X$ and $Z$ of a
generalized CP transformation, proving the reverse implication.

A single symmetric Yukawa matrix within one compatible flavor block always
passes this test: its Takagi factorization makes it diagonal and nonnegative.
Several matrices must admit one common field basis.  Once the scalar CP basis
has been fixed, a useful fermion-basis-invariant obstruction is an imaginary
part such as
\begin{equation*}
  \operatorname{Im}\Tr\left(
    Y^A Y^{B\dagger}Y^C Y^{D\dagger}
  \right),
\end{equation*}
which vanishes whenever all the matrices are real.

For a minimal simultaneous-reality check, take one Weyl bilinear coupled to
two fixed CP-even real scalar directions with coefficients $y_1$ and $y_2$.
A fermion rephasing multiplies both coefficients by the same phase.  A common
real basis therefore exists exactly when
\begin{equation*}
  \operatorname{Im}(y_1y_2^*)=0.
\end{equation*}
This example detects the relative phase that would be missed by making each
coupling real in a separate basis.

This proof uses the involutive CP action assumed in the exercise and basis
changes compatible with the gauge representations.  A pure-gauge term
$\theta F_{\mu\nu}\widetilde F^{\mu\nu}$ is CP odd.  An anomalous chiral
redefinition used in making Yukawa matrices real can shift $\theta$.
Consequently, the complete theory is CP-invariant when the Yukawa couplings
are real in the same basis in which all CP-odd gauge coefficients vanish.
This is the refinement Banks states immediately after the Yukawa criterion.

\BanksImplicitSolution{I-CH05-004}
For one generator every polynomial has the form $f(\eta)=a+b\eta$.  Let $I$
denote a linear integral and translate by an independent odd constant $\xi$.
Translation invariance says
\begin{equation*}
  I(a+b\eta+b\xi)=I(a+b\eta).
\end{equation*}
Since $b$ and $\xi$ are arbitrary, this equality forces $I(1)=0$.  The
normalization $I(\eta)=1$ then gives
\begin{equation*}
  \int\dd\eta\,(a+b\eta)=b.
\end{equation*}
Thus a one-variable Berezin integral equals a left Grassmann derivative.

The left derivative obeys the graded product rule
\begin{equation*}
  \partial_i^L(AB)
  =(\partial_i^L A)B+(-1)^{|A|}A(\partial_i^L B).
\end{equation*}
With the measure convention in the exercise, the iterated integral is
\begin{equation*}
  \int\dd^n\eta\,f
  =\partial_n^L\cdots\partial_1^L f.
\end{equation*}
Expand the polynomial in ordered monomials,
\begin{equation*}
  f(\eta)=\sum_{I=\{i_1<\cdots<i_r\}}
  f_I\eta_{i_1}\cdots\eta_{i_r}.
\end{equation*}
Every monomial of degree below $n$ lacks some generator, so the corresponding
derivative kills it.  The top monomial gives
\begin{equation*}
  \partial_n^L\cdots\partial_1^L
  (\eta_1\cdots\eta_n)=1.
\end{equation*}
It follows that
\begin{equation*}
  \int\dd\eta_n\cdots\dd\eta_1\,f(\eta)
  =f_{\{1,\ldots,n\}}.
\end{equation*}

Anticommutation determines every ordering sign.  For a permutation
$\pi\in S_n$,
\begin{equation*}
  \int\dd\eta_n\cdots\dd\eta_1\,
  \eta_{\pi(1)}\cdots\eta_{\pi(n)}
  =\operatorname{sgn}(\pi).
\end{equation*}
The left derivatives anticommute, so exchanging two neighboring
differentials also multiplies the integral by $-1$.  In two variables,
\begin{align*}
  f&=a+b_1\eta_1+b_2\eta_2+c\eta_1\eta_2,\\
  \int\dd\eta_2\dd\eta_1\,f&=c,\\
  \int\dd\eta_1\dd\eta_2\,f&=-c.
\end{align*}
This sign check fixes the meaning of $\dd^n\eta$ in Banks's coefficient rule.

As a further consistency check, if $\eta=A\xi$, then
$\eta_1\cdots\eta_n=(\det A)\xi_1\cdots\xi_n$.  Coefficient extraction
therefore requires the Grassmann Jacobian
\begin{equation*}
  \dd^n\eta=(\det A)^{-1}\dd^n\xi.
\end{equation*}
The determinant and Pfaffian formulas that follow in Banks's discussion arise
by extracting the top monomial of the corresponding Grassmann Gaussian.

\end{bankssolutions}
